\section{The allocation problem}
\label{sec:problem}

\subsection{Objective}

Let $\mathbf x_{\mathrm{ref}}(t)$ be a reference trajectory on $[0,T]$ and let

\begin{equation}
  \Delta\mathbf a(t,N) \;=\;
  \mathbf a_N\bigl(\mathbf x_{\mathrm{ref}}(t)\bigr)
  -\mathbf a_{\mathrm{ref}}\bigl(\mathbf x_{\mathrm{ref}}(t)\bigr)
  \label{eq:defect}
\end{equation}

be the truncation acceleration defect of degree $N$ measured along it. To first
order the induced \emph{state} perturbation is
$\delta\mathbf x(t)=\int_0^t \Phi(t,\tau)\,\mathbf B\,\Delta\mathbf a(\tau)\,\dd\tau
\in\mathbb{R}^6$, with $\Phi$ the reference state-transition matrix and
$\mathbf B$ mapping acceleration into the velocity block
\citep[Ch.~4]{tapley2004statistical}. The metric this paper reports is a
\emph{position} RMS, so the position block must be selected before the norm is
taken. With $\mathbf H_r=[\,\mathbf I_3\;\;\mathbf 0\,]$,

\begin{equation}
  E^2 \;=\; \frac{1}{T}\int_0^{T}
  \Bigl\lVert \mathbf H_r\!\int_0^{t}\Phi(t,\tau)\,\mathbf B\,
  \Delta\mathbf a(\tau)\,\dd\tau \Bigr\rVert^2 \dd t .
  \label{eq:objective}
\end{equation}

Taking the norm of the full six-vector would mix position and velocity units
and would not be the quantity the campaign measures; $\mathbf H_r$ is carried
explicitly through every expression below for that reason.

The decision variable is the degree history $N(\cdot)$, and the cost is the
quadratic work $\int N(t)^2\,\dd t$ it implies.

\subsection{Why the separable treatments do not apply}

Equation~\eqref{eq:objective} takes the norm \emph{after} the integral, and
that is the whole of the difficulty. Once the objective is discretized, Section~\ref{sec:problem-discretization}
puts it in the form $\mathbf u^{\top}\mathbf Q\,\mathbf u$, and the three
functionals separating this paper from its predecessor are then literally three
readings of \emph{one} matrix:

\begin{equation}
  \underbrace{\textstyle\sum_i \lVert\Delta\mathbf a_i\rVert^2\Delta t_i}
             _{\text{(i) defect norm: no }\mathbf Q}
  \qquad
  \underbrace{\textstyle\sum_i \Delta t_i\,
     \Delta\mathbf a_i^{\top}\mathbf K_i\,\Delta\mathbf a_i}
             _{\text{(ii) diagonal kernel}}
  \qquad
  \underbrace{\mathbf u^{\top}\mathbf Q\,\mathbf u}
             _{\text{(iii) all of }\mathbf Q}
  \label{eq:three-functionals}
\end{equation}

Functionals (i) and (ii) are separable in the per-epoch degrees --- (ii)
because a block-diagonal quadratic form decouples --- so a single multiplier
decides each epoch independently \citep{everett1963generalized}. Functional
(iii) is not: changing $N_i$ changes the sum against which every other epoch's
contribution is measured.

Writing the ladder this way removes an arbitrary choice: the sensitivity weight
in (ii) is not borrowed from anywhere, it is the objective's own coupling
evaluated on its diagonal. \emph{RQ2 is therefore exactly: how much of the gain
lives off the diagonal?} The previous paper's finding --- that a signed,
remaining-horizon-weighted statistic still ranks two policies the wrong way
round on \ph{n} of \ph{n} orbits --- says the answer is not zero.

\paragraph{The diagonal restriction has to be taken in the kernel, not in the
matrix.} Written continuously, $J$ is a double integral against a kernel,
$J=\int\!\!\int\Delta\mathbf a(\tau)^{\top}\mathbf K(\tau,\sigma)
\Delta\mathbf a(\sigma)\,\dd\tau\,\dd\sigma$, and the set $\tau=\sigma$ has
measure zero. ``Dropping the off-diagonal blocks'' of the discrete
$\mathbf Q$ is therefore not a definition one can take literally: since
$\mathbf u_i\propto\Delta t_i$, the raw diagonal term
$\mathbf u_i^{\top}\mathbf Q_{ii}\mathbf u_i$ carries $\Delta t_i^2$ and shrinks
as the grid is refined, so a criterion built on it would depend on the
accumulation resolution --- which, with the refined grid of
Section~\ref{sec:problem-grids}, is not even uniform along the arc.

What is well defined is the kernel \emph{on} the diagonal, a local sensitivity
$\mathbf K(t,t)$:

\begin{equation}
  \mathbf K(t,t) \;=\; \frac{1}{T}\int_t^T
  \mathbf B^{\top}\Phi(\tau,t)^{\top}\mathbf H_r^{\top}\mathbf H_r\,
  \Phi(\tau,t)\,\mathbf B\;\dd\tau ,
  \qquad
  \mathbf K_i \;=\; \frac{1}{T}\,\mathbf B^{\top}\Phi(t_0,m_i)^{\top}
  \mathbf A_i\,\Phi(t_0,m_i)\,\mathbf B ,
  \label{eq:local-sensitivity}
\end{equation}

the discrete form following from Eq.~\eqref{eq:Q-structure} and the cocycle
property. Functional (ii) is then $\sum_i\Delta t_i\,
\Delta\mathbf a_i^{\top}\mathbf K_i\Delta\mathbf a_i$, one power of $\Delta t$
per term like (i), and it converges to a physical limit under refinement.
Equivalently it is $\sum_i \mathbf u_i^{\top}\mathbf Q_{ii}\mathbf u_i/\Delta t_i$
--- the raw diagonal with the extra $\Delta t$ removed.

Nothing else moves. $\mathbf K_i$ is built from the same suffix sequence
$\mathbf A_i$, so it costs nothing extra; and (i) and (iii) were already
grid-consistent, (iii) because a double sum over $\mathbf u$ carries
$\Delta t_i\Delta t_k$ and is a genuine double integral. The derivation of
$\mathbf Q$, the prefix--suffix structure and the Frank--Wolfe certificate are
untouched.

\subsection{Discretization}
\label{sec:problem-discretization}

Fix an accumulation grid of $M$ cells with \emph{edges}
$e_0=0<e_1<\dots<e_M=T$ and widths $\Delta t_i=e_{i+1}-e_i$ --- not assumed
uniform, for the reason given in Section~\ref{sec:problem-grids}. Two index
sets are in play and they are deliberately distinct: $i=0,\dots,M-1$ runs over
\emph{cells}, whose midpoints $m_i$ carry the defect, and $j=0,\dots,M$ runs
over \emph{edges}, where the objective is sampled. Keeping them apart is what
the next paragraph is about; conflating them puts a term in the coupling that
violates causality.

With $\mathbf H_r=[\,\mathbf I_3\;\;\mathbf 0\,]$ selecting the position block,

\begin{equation}
  \mathbf u_i(N) = \Phi(t_0,m_i)\,\mathbf B\,\Delta\mathbf a(m_i,N)\,
  \Delta t_i ,
  \qquad
  \mathbf S_j = \sum_{i\le j-1}\mathbf u_i,
  \qquad
  \delta\mathbf r(e_j) = \mathbf H_r\,\Phi(e_j,t_0)\,\mathbf S_j ,
  \label{eq:u-and-S}
\end{equation}

using only the cocycle property $\Phi(e_j,m_i)=\Phi(e_j,t_0)\Phi(t_0,m_i)$.
The prefix runs to $j-1$ because cell $i$ occupies $[e_i,e_{i+1}]$ and is
therefore felt at edge $e_j$ exactly when $j\ge i+1$: at its own left edge the
impulse has not yet happened, and $\mathbf S_0=\mathbf 0$.

\paragraph{Which quadrature rule the prefix sum permits.} That
$\mathbf S_j$ is a \emph{plain} prefix sum is not a notational convenience; it
requires the inner weight of epoch $i$ to be independent of the outer epoch
$j$, and that requirement selects the rule. A trapezoid rule for the inner
integral has endpoint coefficient $(t_j-t_{j-1})/2$, which depends on $j$, so
carrying it would replace Eq.~\eqref{eq:Q-structure} by that structure plus a
$j$-dependent correction and forfeit the $\max(i,k)$ form. A rectangle rule
keeps the structure but is first-order, which at two samples per correlation
time is not a rounding detail.

We therefore place the accumulation nodes at the \emph{midpoints} of the grid
cells and sample the outer integral at the cell \emph{edges}. The inner weight
is then the full cell width --- independent of $j$, so the $\max(i,k)$
structure survives --- while the rule itself is the midpoint rule and
second-order, and $\mathbf S_j$ is the exact prefix sum of that
discretization through edge $j$. The outer integral carries no prefix sum and
is free to use the trapezoid rule on the edges. Both weight sets are
non-negative, which is what Section~\ref{sec:problem-structure} needs.

The price of separating the two grids is that the suffix in
Eq.~\eqref{eq:Q-structure} starts one index later than it would on a single
grid, and that offset is not cosmetic: it is the difference between a
discretization in which a cell influences the state \emph{after} it and one in
which it also influences the state at its own left edge.

The outer integral of Eq.~\eqref{eq:objective} is likewise a quadrature, with
weights $\omega_j\ge0$ summing to $T$, so the discretized objective is a
quadratic form in the stacked vector
$\mathbf u=(\mathbf u_0,\dots,\mathbf u_{M-1})$,

\begin{equation}
  J(\mathbf u) \;=\; \frac{1}{T}\sum_{j} \omega_j\,
  \mathbf S_j^{\top}\mathbf M_j\mathbf S_j
  \;=\; \mathbf u^{\top}\mathbf Q\,\mathbf u ,
  \qquad
  \mathbf M_j \;=\; \Phi(t_j,t_0)^{\top}\mathbf H_r^{\top}\mathbf H_r\,
  \Phi(t_j,t_0) ,
  \label{eq:quadratic-form}
\end{equation}

so that each $\mathbf M_j$ is a symmetric positive-semidefinite $6\times6$
block and no dimensional ambiguity is left in the notation. On a uniform grid
$\omega_j=\Delta t$ and $T=M\Delta t$, and Eq.~\eqref{eq:quadratic-form} reduces to
the unweighted average $M^{-1}\sum_j\mathbf S_j^{\top}\mathbf M_j\mathbf S_j$;
every expression below is written in the weighted form so that the two cases
are one derivation rather than two.

\paragraph{The decision variable lives on the decision grid.} The sum in
Eq.~\eqref{eq:quadratic-form} runs over accumulation epochs, but the degree
cannot be chosen $M$ times. Let $g(i)=q$ map an accumulation epoch to its
decision interval, $I_q=\{i: g(i)=q\}$, and write

\begin{equation}
  N_i = N_{g(i)} ,
  \qquad
  W_q \;=\; \sum_{i\in I_q}\Delta t_i
  \label{eq:interval-weight}
\end{equation}

for the interval's total time weight. The allocation problem is then

\begin{equation}
  \min_{(N_1,\dots,N_{K_{\mathrm{dec}}}) \,\in\, \mathcal{N}^{K_{\mathrm{dec}}}}
  \; \mathbf u^{\top}\mathbf Q\,\mathbf u
  \qquad \text{subject to} \qquad
  \sum_q W_q\,N_q^2 \;\le\; B \;=\; B_1 T ,
  \qquad K_{\mathrm{dec}} \ll M .
  \label{eq:allocation}
\end{equation}

The budget is a time integral of the instantaneous work rate, not a count of
samples: $B_1$ of Section~\ref{sec:setup-budgets} is the time average
$\langle N^2\rangle_t$, so the total over an arc of length $T$ is $B_1T$. Weighting by $W_q$ rather than by
the sample count $\lvert I_q\rvert$ is what keeps this correct when the
accumulation grid is refined near perilune --- an unweighted count would charge
the budget as though perilune occupied most of the arc simply because it is
sampled most densely. On a uniform grid $W_q=\lvert I_q\rvert\Delta t$ and the
two agree.

This is not only an implementability constraint. Optimizing over $M$ variables
would give the benchmark more \emph{switching freedom} than any deployable
controller has, and its advantage would then be partly a statement about grid
resolution rather than about trajectory-aware information. Every allocation
compared in this paper --- benchmark and controller alike --- is piecewise
constant on the same decision grid.

The result is an integer quadratic program with one knapsack constraint. Its
structure, however, is unusually favourable.

\subsection{The structure that makes it cheap}
\label{sec:problem-structure}

Because $\mathbf S_j$ is a prefix sum, the block of $\mathbf Q$ coupling cells
$i$ and $k$ collects exactly those edges that see both. Cell $i$ is seen at
edge $j$ when $j\ge i+1$, so the pair is seen when $j\ge\max(i,k)+1$:

\begin{equation}
  \mathbf Q_{ik} \;=\; \frac{1}{T}\!\!\sum_{j\ge\max(i,k)+1}\!\!\omega_j\,
  \mathbf M_j
  \;=\; \frac{1}{T}\,\mathbf A_{\max(i,k)} ,
  \qquad
  \mathbf A_i = \sum_{j= i+1}^{M} \omega_j\,\mathbf M_j .
  \label{eq:Q-structure}
\end{equation}

The offset by one is the whole content of separating cells from edges, and it
is worth stating why it cannot be dropped. Writing $\mathbf A_i=\sum_{j\ge i}$
instead would add $\omega_i\mathbf M_i$ to every suffix block, which credits
cell $i$ with displacing the state at $e_i$ --- before its own impulse. The
resulting form is not a small perturbation of the right one: on a random
five-cell instance it misstates $J$ by $26\%$, and it breaks the identity
between $\mathbf K_i$ and the diagonal of the kernel.

The quadrature weights enter only inside the suffix sequence, so the structure
that makes the problem tractable is untouched by a non-uniform grid. They are
required to be non-negative, which the trapezoid rule used here satisfies:
with $\omega_j\ge0$ and each
$\mathbf M_j$ positive semidefinite, $\mathbf Q$ is positive semidefinite, and
that is what the convex relaxation of Section~\ref{sec:algorithm-certificate}
rests on.

$\mathbf Q$ therefore has no free entries at all: it is generated by the single
suffix sequence $\mathbf A_i$, whose members are symmetric positive-semidefinite
$6\times6$ blocks --- the same size as $\mathbf M_j$, since $\mathbf u_i$ is a
state perturbation and the position selection happens inside $\mathbf M_j$
rather than on $\mathbf u_i$. The cross term needed by any coordinate method
follows from one prefix sum and one suffix sum,

\begin{equation}
  \mathbf c_i \;:=\; \sum_{k}\mathbf Q_{ik}\mathbf u_k
  \;=\; \frac{1}{T}\Bigl[\;
    \mathbf A_i \sum_{k\le i}\mathbf u_k
    \;+\; \sum_{k>i}\mathbf A_k\,\mathbf u_k
  \;\Bigr] ,
  \label{eq:cross-term}
\end{equation}

so evaluating $\mathbf c_i$ at \emph{all} $M$ epochs costs $\mathcal{O}(M)$
rather than the $\mathcal{O}(M^2)$ a dense $\mathbf Q$ would require, and the
gradient $\nabla J = 2\mathbf Q\mathbf u$ is available at the same cost. A
coordinate \emph{sweep} additionally tries $\lvert\mathcal{N}\rvert$ candidates
per decision interval and so costs
$\mathcal{O}(M\lvert\mathcal{N}\rvert)$; the two figures are quoted separately
throughout because they measure different things.

Grouping the epochs into decision intervals does not spoil this. For
$I_q=\{i_1<\dots<i_m\}$ the within-block term telescopes, because
$\max(i_p,i_{p'})=i_{p}$ whenever $p'<p$:

\begin{equation}
  \sum_{i,k\in I_q}\mathbf u_i^{\top}\mathbf Q_{ik}\mathbf u_k
  \;=\;
  \frac{1}{T}\Biggl[
  \sum_{p}\mathbf u_{i_p}^{\top}\mathbf A_{i_p}\mathbf u_{i_p}
  \;+\; 2\sum_{p}\Bigl(\sum_{p'<p}\mathbf u_{i_{p'}}\Bigr)^{\!\top}
  \mathbf A_{i_p}\mathbf u_{i_p}\Biggr] ,
  \label{eq:within-block}
\end{equation}

which one running prefix sum evaluates in $\mathcal{O}(m)$, and the
between-block term is Eq.~\eqref{eq:cross-term} with the in-block part
subtracted. A full sweep over all $K_{\mathrm{dec}}$ decision variables and all
candidate degrees therefore remains
$\mathcal{O}(M\lvert\mathcal{N}\rvert)$. Section~\ref{sec:algorithm} uses this.

The terminal-position objective is the special case in which the outer
quadrature collapses onto the final edge, $\omega_j = T\,\delta_{jM}$, so that
$\mathbf A_i = T\,\mathbf M_M$ for every cell and

\begin{equation}
  J_T \;=\; \bigl\lVert\,\mathbf H_r\,\Phi(e_M,t_0)\textstyle\sum_i\mathbf u_i
  \,\bigr\rVert^2 ,
  \label{eq:terminal}
\end{equation}

with the transport matrix retained because the $\mathbf u_i$ have been carried
back to $t_0$; dropping it would score the accumulated perturbation in the
wrong frame. Everything below covers this case unchanged.

\subsection{Two grids, not one}
\label{sec:problem-grids}

The decision variable and the integrand live on different time scales, and
conflating them is a measurement error rather than a modelling choice.

The degree schedule $N(t)$ varies on the orbital time scale. The defect
$\Delta\mathbf a$ does not: its direction is set by the leading omitted
harmonics, whose half-wavelength at the spacecraft's radius is $\pi r/N$, so
along track it decorrelates over
\begin{equation}
  \tau_{\mathrm{corr}} \;\approx\; \frac{\pi r}{N\,v} .
  \label{eq:decorrelation}
\end{equation}
A signed integral sampled coarser than that is aliased, even where the
corresponding magnitude statistic is converged --- and the previous paper
verified convergence only of the magnitude, which moved by \ph{pct} between a
120~s and a 10~s cadence.

The speed in Eq.~\eqref{eq:decorrelation} is the inertial $v$, and that is a
deliberately conservative proxy. What carries the vehicle across the field's
texture is the body-fixed transverse rate $r\lVert\dot{\hat{\mathbf r}}_{\mathrm{BF}}\rVert$,
which equals $v$ only where the motion is transverse; on an eccentric arc the
radial component inflates $v$ without moving the ground track, so
Eq.~\eqref{eq:decorrelation} runs short and the grid and probe are costed
high rather than low. The convergence sweep of Section~\ref{sec:campaign}
settles the question empirically --- it varies the sampling and asks where the
schedule stops moving --- so nothing rests on the proxy being tight.

Equation~\eqref{eq:decorrelation} is not one number. Both $r$ and $v$ swing by
factors of several on the eccentric orbits studied here, and $N$ with them, so
$\tau_{\mathrm{corr}}$ runs from \ph{value}~s near a 50~km perilune at $N=300$
to \ph{value}~s near apolune, a range of about \ph{value}. A uniform
accumulation grid fine enough for perilune is therefore mostly spent where
nothing happens, and one coarse enough for apolune aliases the passage that
matters. The grid is refined on $\tau_{\mathrm{corr}}$ instead, which also
requires the budget to be accumulated with time weights,
$B_1 = \sum_i N_{g(i)}^2\,\Delta t_i / T$; on a uniform grid this reduces to
the mean of Section~\ref{sec:setup-budgets}. \wpref{WP1, WP4}

We therefore separate

\begin{itemize}
  \item the \emph{accumulation grid} $\Delta t_{\mathrm{acc}}$, on which
    $\Delta\mathbf a$, $\Phi$ and the partial sums are formed, chosen fine
    enough that $J$ and the recovered schedule are converged
    (Section~\ref{sec:setup-grids}); and
  \item the \emph{decision grid} $\Delta t_{\mathrm{dec}}$, on which $N$ is held
    constant, chosen coarse enough to be implementable.
\end{itemize}

A schedule is thus piecewise constant on the decision grid while its
consequences are integrated on the accumulation grid. Section~\ref{sec:setup-grids}
reports the convergence of both, and Section~\ref{sec:ablations} reports what
coarsening the decision grid costs. \wpref{WP4}
